\section{Risk Communication and the Verification Asymmetry}
\label{sec:communication}

The preceding section assumes the actor can identify exercise pathways from
its own world model. In a multi-agent population, risk information also flows
through communication: agents report risk assessments to each other. This
section analyzes how risk claims enter the trust model and the fundamental
epistemological challenge they create.

\subsection{Risk Claims as Counterfactual Predictions}

A risk claim from agent $a_j$ has the form: ``capability tuple $(d_1, \ldots,
d_m)$ held by agent $a_k$ has exercise pathway $\Pathway$ with contraction
magnitude $|\Delta \volL|$ and probability $p$.'' This enters the GFM actor's
world model $\WorldModel$ as a communicated belief about future states.

Risk claims differ from capability claims in a fundamental way. A capability
claim (``I can run a marathon'') is verifiable through the observation
channel: watch the agent run. A risk claim (``this capability combination is
dangerous'') is a prediction about a counterfactual future: what would happen
if the capabilities were exercised together. The prediction is not directly
verifiable until the exercise either happens or does not. This
distinction motivates a separate trust channel.

The foundational paper's trust factor $T_j$ measures \emph{prediction
consistency}: how well the model $M_j(G)$ tracks $a_j$'s observed behavior
\citep{lasser2026gfm}. A predictably destructive agent can maintain high
$T_j$. Using $T_j$ directly to weight risk claims would allow a
high-consistency agent to have outsized influence on risk assessments before
any risk-specific evaluation occurs. We therefore introduce a separate
\emph{risk-trust} factor $T_j^{\mathrm{risk}} \in [0, 1]$, initialized to
a neutral prior and updated through the causal-counterfactual evaluation
described below.

\begin{definition}[Risk Claim]
\label{def:risk_claim}
A \emph{risk claim} from agent $a_j$ is a tuple $(S, \Pathway, |\Delta
    \volL|, p)$ where $S = \{d_1, \ldots, d_m\}$ is a capability set,
    $\Pathway$ is an exercise pathway, $|\Delta \volL|$ is the claimed
    contraction magnitude, and $p$ is the claimed exercise probability. The
    claim enters the actor's risk model weighted by risk-trust:
\begin{equation}
\label{eq:risk_claim_weight}
\Risk_j(S; a_k, t) = T_j^{\mathrm{risk}} \cdot p \cdot |\Delta \volL|
    \cdot \mathbb{E}[\gamma^{T_\Pathway}]
\end{equation}
where the expectation is over the stochastic execution time $T_\Pathway$.
\end{definition}

Multiple agents may report different risk assessments for the same
capability tuple. The actor aggregates them:
\begin{equation}
\label{eq:risk_aggregate}
\Risk_{\mathrm{agg}}(S; a_k, t) = \max\!\left(
    \Risk_{\mathrm{self}}(S; a_k, t),\;
    \max_j \Risk_j(S; a_k, t)
\right)
\end{equation}
taking the maximum of the actor's own assessment and all communicated
assessments. The $\max$ operator is conservative: a single high-trust agent
reporting high risk is sufficient to trigger evaluation of preemptive
restriction (Proposition~\ref{prop:preemptive}). This is appropriate for
catastrophic risks where false negatives (missing a real threat) are far
more costly than false positives (investigating a non-threat).

\subsection{The Verification Asymmetry}

\begin{remark}[The Verification Asymmetry]
\label{rem:verification}
Risk-trust convergence has a structural asymmetry absent from
    capability-trust convergence. Capability claims converge because the
    observation channel provides ground truth: the agent either demonstrates
    the capability or does not, and the prediction residual $r_k(t)$
    accumulates accordingly. Risk claims have no comparable ground truth
    mechanism:
\begin{enumerate}
\item If the risk materializes (the exercise pathway is executed), the
    risk claim is confirmed---but the catastrophe has already occurred, and
    the confirmation is too late to be useful for prevention.
\item If the risk is \emph{prevented} (the actor restricts the keystone
    capability, blocking the exercise pathway), the risk claim is neither
    confirmed nor falsified. The counterfactual (``what would have happened
    without restriction?'') is never observed.
\item If the risk does not materialize and no prevention was taken, the
    risk claim appears to be falsified---but the non-materialization may
    reflect the agent's choice not to exercise the pathway rather than the
    pathway's non-existence.
\end{enumerate}
\end{remark}

Case~2 is the core problem, which we term the \emph{Wamura problem} after the
example developed in \citet{lasser2026poset}: a successful risk intervention
removes the very evidence that would have justified it, and the convergence
timescale for risk-trust can exceed the timescale on which trust normally
accumulates. The Wamura problem is not a deficiency of this framework; it is
an inherent feature of risk reasoning that any trust model must confront.

\subsection{Causal-Counterfactual Resolution}

The verification asymmetry cannot be fully resolved within an observational
trust model. We propose a partial resolution through causal counterfactual
evaluation \citep{pearl2009causality, spirtes2000causation}.

\paragraph{Causal risk evaluation.}
\label{def:causal_risk}
Given a risk claim $(S, \Pathway, |\Delta \volL|, p)$ from agent $a_j$, the
actor evaluates the claim by constructing the causal chain from capability
possession to exercise, in four steps:
\begin{enumerate}
\item \textbf{Prerequisite verification.} Does agent $a_k$ actually possess
    all capabilities in $S$? (Verifiable through the observation channel.)
\item \textbf{Pathway plausibility.} Does the actor's own world model
    $\WorldModel$ confirm that the exercise pathway $\Pathway$ exists? (The
    actor simulates the pathway under its own causal model, independent of
    $a_j$'s claim.)
\item \textbf{Probability calibration.} Is $a_j$'s claimed probability $p$
    consistent with the actor's own estimate based on $T_k$ and $a_k$'s
    behavioral history?
\item \textbf{Magnitude plausibility.} Is the claimed contraction magnitude
    $|\Delta \volL|$ consistent with the number of agents and capabilities
    affected by the pathway?
\end{enumerate}

The causal evaluation does not require the catastrophe to happen. It requires
the actor to independently verify the \emph{structure} of the risk claim:
Does the capability combination exist? Does the exercise pathway exist in the
actor's own model? Is the claimed probability in the right range? Each of
these is evaluable without observing the outcome.

\paragraph{Trust update for risk claims.} The actor updates
$T_j^{\mathrm{risk}}$ based on the structural plausibility of the claim
rather than the outcome:
\begin{itemize}
\item If the claim's structure is independently verified (the pathway exists
  in $\WorldModel$, the probability estimate is calibrated, the magnitude is
  plausible), $T_j^{\mathrm{risk}}$ increases regardless of whether the
  catastrophe occurs.
\item If the claim's structure is not verified (the pathway does not exist in
  $\WorldModel$, or the probability is implausibly high),
  $T_j^{\mathrm{risk}}$ decreases.
\item If the catastrophe occurs, all agents who reported the risk receive a
  large $T_j^{\mathrm{risk}}$ boost; agents who had the information to
  report but did not receive a penalty.
\end{itemize}

Unlike the behavioral trust model in the foundational paper (which has
explicit prediction-residual dynamics), the $T_j^{\mathrm{risk}}$ update
is described qualitatively. A formal update equation analogous to the
foundational paper's EWMA dynamics would strengthen the risk-communication
mechanism but is deferred to future work.

This does not fully resolve the Wamura problem---a risk claim whose pathway is
plausible but whose probability is genuinely uncertain remains in a trust
limbo until either the catastrophe occurs or enough time passes to revise the
probability estimate downward. But it narrows the asymmetry: risk-trust
converges on the \emph{quality of the causal model}, not on the
\emph{occurrence of the outcome}.
